
\documentclass[final,ignorenonframetext,compress]{beamer}

\mode<beamer>
{
\usetheme{Bonn}
\pgfdeclareimage[height=0.8cm]{university-logo}{images/Logo_UBo_h24_4c-crop}
\setbeamertemplate{footline}[default]
\logo{\pgfuseimage{university-logo}}
\setbeamertemplate{navigation symbols}{}
}
\mode<handout>{
\usetheme{default}
\usecolortheme{dove}
\setbeamercolor{background canvas}{bg=black!5}
\usepackage{pgfpages}
\pgfpagesuselayout{8 on 1}[a4paper,border shrink=5mm]
}


\DeclareMathOperator*{\mymax}{max}
\DeclareMathOperator*{\argmax}{arg\,max}
\DeclareMathOperator*{\argmin}{arg\,min}
\DeclareMathOperator*{\ForAll}{\bigwedge}
\DeclareMathOperator*{\Exists}{\bigvee}
\newcommand{\W}{\mathcal{W}}

\newcommand*\kreis[2]{\unitlength0.7ex\begin{picture}(2.5,2.5)%


	\color{#2}%
	\put(0.75,0.75){\circle*{3}}%
	\color{black}%
	\put(0.75,0.75){\circle{3}}%
	\put(0.7,0.7){\makebox(0,0){#1}}\end{picture}}
	\newcommand\itplus{\kreis{$+$}{green}}
	\newcommand\itminus{\kreis{$-$}{red}}
	\newcommand\prob[2]{\ensuremath{p_{\mathrm{#1}}\left( #2 \right)}}
	\newcommand\bedprob[3]{\ensuremath{p_{\mathrm{#1}}\left( #2 \left|#3\right. \right)}}

    \newcommand\etal{et.\,al.}

	\usepackage[english]{babel}
	% avant, courier, chancery, times, palatino, bookman, 
	% newcent, utopia, charter
	%\usepackage{times}

	\usepackage[latin1]{inputenc}
	\usepackage[T1]{fontenc}
	\usepackage{multicol}
	\usepackage{mdwtab}
	\usepackage{booktabs}
    \usepackage{tabularx}
	\usepackage{txfonts,textcomp}
	\usepackage{amsmath}
	\usepackage{amsfonts}
    \usepackage{natbib}

	\setlength\parindent{0cm}
	\setlength\parskip{2mm}
	\setbeamerfont{gross}{size=\large}
	\setbeamerfont{quetsch}{size=\tiny}
	\setbeamerfont{klein}{size=\footnotesize}
	\newcommand\w[1]{\mathbf{#1}}

	% example list spacing
	\newcommand\ml{\setlength\itemindent{-11mm}}

    %\newcommand{\newblock}{}


	\title{
Learning Depth-Sensitive Conditional Random Fields for Semantic Segmentation of RGB-D Images
    } 
    \author[short]{Andreas M\"uller}%
	\institute[UBO]{%
    Advisor: Sven Behnke\\
    Institute for Autonomous Intelligent Systems\\
	Rheinische Friedrich-Wilhelms-Universit\"at\\[4mm]
	\includegraphics[width=.2\linewidth]{images/Logo_UBo_h24_4c-crop}\\[3mm]
	Bonn, Germany}


	\begin{document}

	\begin{frame}[plain]
		\titlepage
	\end{frame}

    \section*{Introduction}

	\begin{frame}
		\frametitle{Semantic Segmentation of Structure Classes}
        %FIXME
            \begin{figure}
                \begin{visibleenv}<1-2>
                \includegraphics[width=.4 \textwidth]{images/00062_image}
                \end{visibleenv}
                \begin{visibleenv}<3>
                \includegraphics[width=.4 \textwidth]{images/00062_pipeline_sp}
                \end{visibleenv}
                \begin{visibleenv}<2->
                    \includegraphics[width=.4 \textwidth]{images/00062_gt}
                \end{visibleenv}
            \end{figure}

	\end{frame}

    \section{Structured prediction formulation}

    \begin{frame}
        \frametitle{Semantic Segmentation as Structured Prediction}
        %FIXME
        %\begin{figure}
            %\includegraphics<1-2>[width=.6\linewidth]{images/scene_sp_gt}
            %\includegraphics<3->[width=.6\linewidth]{images/scene_sp_graph}
        %\end{figure}
        \vspace{-5mm}
        \begin{itemize}
            \item<1->Predict class for each superpixel.
            \item<2->Predictions clearly not independent.
            \item<3->Use neighborhood structure to model correlations.
        \end{itemize}
    \end{frame}

    \begin{frame}
        \frametitle{Structured Prediction}
        Learn prediction function of the form
        \[f(x, w) := \argmax_{y \in \mathcal{Y}}  w^T \psi(x, y) \]
        \begin{visibleenv}<2->
            Pairwise model:
            \[w^T \psi(x, y) = \sum_{(i, j) \in E} w_{i,j} \psi_{i,j}(x, y_i, y_j) + \sum_i w_i \psi_i(x, y_i)\]
        \end{visibleenv}
        Objective
        % FIXME
    \end{frame}

    \section{Model and Features}
    \begin{frame}
        \frametitle{Superpixels}
        % image + formular!
    \end{frame}

    \begin{frame}
        \frametitle{Random Forest Unary Potentials}
        %TODO reference
    \end{frame}

    \begin{frame}
        \frametitle{Absolute height potentials}
        \begin{center}
            \includegraphics[width=.6\linewidth]{images/height_success}\\
            \vspace{3mm}
            \includegraphics[width=.6\linewidth]{images/height_failure}
        \end{center}
        \begin{visibleenv}<2->
            Formular %TODO
        \end{visibleenv}

    \end{frame}


    \begin{frame}
        \frametitle{Color Contrast}
        \begin{center}
            \includegraphics[width=.8\linewidth]{images/00062_pipline_feature_pairwise_1}
        \end{center}
        \begin{visibleenv}<2->
            Formular %TODO
        \end{visibleenv}

    \end{frame}

    \begin{frame}
        \frametitle{Vertical Alignment}
        \begin{center}
            \includegraphics[width=.8\linewidth]{images/00062_pipline_feature_pairwise_4}
        \end{center}
        \begin{visibleenv}<2->
            Formular %TODO
        \end{visibleenv}
    \end{frame}

    \begin{frame}
        \frametitle{Depth Difference}
        \begin{center}
            \includegraphics[width=.8\linewidth]{images/00062_pipline_feature_pairwise_2}
        \end{center}
        \begin{visibleenv}<2->
            Formular %TODO
        \end{visibleenv}
    \end{frame}
    

    \begin{frame}
        \frametitle{Pairwise Normal Orientation}
        \begin{center}
            \includegraphics[width=\linewidth]{images/normal_feature}
        \end{center}
        \begin{visibleenv}<2->
            Formular %TODO
        \end{visibleenv}
    \end{frame}


    \begin{frame}
        \frametitle{Summary of Pairwise Potentials}
            \noindent\emph{$\bullet$~Constant.}
                A constant feature allows to model general neighborhood relations.


            \noindent\emph{$\bullet$~Color Contrast.}
                We employ a non-linear color contrast, as is common in the computer
                vision literature, between the superpixel mean colors $c_i$ and $c_j$:
                $\exp\left(-\gamma \lVert c_i - c_j \rVert^2\right)$.


            \noindent\emph{$\bullet$~Vertical Alignment.}
                We model the directed angle between superpixel centers in the 2D image
                plane.  This allows the  model to learn that ``structure'' is above
                ``floor'', but not the other way around.


            \noindent\emph{$\bullet$~Depth Difference.}
                We include the signed depth difference between superpixels, which
                allows the model to detect depth discontinuities that are not
                represented in the 2D neighborhood graph of the superpixels.


            \noindent\emph{$\bullet$~Normal Orientations.}
                Differences in normal vector orientation are a strong clue on
                whether two superpixels belong to the same surface, and therefore the
                same structural class.
                We compute the 3D orientation of normals using the method of \citet{holz_2011_robocup},
                as implemented in the point cloud library (pcl)\footnote{\url{http://pointclouds.org}}.
                All normals within a superpixel are then averaged, to get a single orientation for each superpixel.
                The feature is computed as the difference of $\frac{\pi}{4}$ and the (undirected) angle between the normals belonging
                to two adjacent superpixels.  
                The change
                in normal orientation highlights that pillow and wall are distinct
                objects, even though there is no strong distinction in color or depth.
    \end{frame}

    \begin{frame}
        \frametitle{How intractable are loopy models?}
        \begin{itemize}
            \item In general, inference in loopy models is NP-hard.
            \item Learning relies on inference.
            \item Combining several inference algorithms, we can get exact results in practice. %FIXME reference
        \end{itemize}
    \end{frame}

    \begin{frame}
        \frametitle{Summary}
        \begin{center}
            \includegraphics[width=.9\linewidth]{images/teaser}
        \end{center}
    \end{frame}

    \section{Experiments}
    
    \begin{frame}
        \frametitle{Dataset}
        %TODO
    \end{frame}
    

    \begin{frame}
        \frametitle{Experimental Protocol}
        %TODO
    \end{frame}

    \begin{frame}
        \frametitle{Results}
        \begin{table}[t]
            \caption{Quantitative comparison of the proposed method with the
            literature.}
        \begin{tabularx}{\linewidth}{@{\extracolsep{\fill}}lcccccc}
        \toprule
                                & ground        &  structure    & furniture     & props         & class average   & pixel average\\
        \cmidrule(r){1-7}

        RF                              &         90.8  &   81.6        & 67.9          & 19.9          &  65.0        &  68.3 \\
        RF + SP                         &         92.5  &   83.3        & \textbf{73.8} & 13.9          &  65.7        &  70.1 \\ 
        RF + SP + SVM                   &         94.4  &   79.1        & 64.2          & \textbf{44.0} &  70.4        &  70.3 \\
        RF + SP + CRF                   & \textbf{94.9} &   78.9        &          71.1 & 42.7          &\textbf{71.9} &  \textbf{72.3} \\
        \cmidrule(r){1-7}
        \citet{SilbermanECCV12}         &         68    &   59          & 70           & 42            &  59.6        & 58.6 \\
        \citet{couprie-iclr-13}         &         87.3  & \textbf{86.1} & 45.3         & 35.5          &  63.5        & 64.5 \\
        \cmidrule(r){1-7}
        \citet{stueckler2013}$^\dagger$ & \textbf{95.6} &   83.0        & \textbf{75.1}& 14.2          &  67.0        & 70.9 \\

        \bottomrule
        \end{tabularx}
        \quad\\The best value in each column is printed in bold$^\dagger$. The upper part of
        the table shows contributions by different parts of our pipeline. RF stands for random forest prediction, RF + SP for aggregated
        random forests prediction within superpixels, RF + SP + SVM for an SVM trained on the unary potentials, and RF + SP + CRF is
        our proposed pipeline. We optimized our approach for class average
        accuracy.\\
        $^\dagger$ \footnotesize Note that the work of \citet{stueckler2013} is not directly
        comparable, as they integrated information over multiple frames, and did not
        measure accuracy for pixels without valid depth measurement.
        \vspace{-5mm}
        \end{table}
    \end{frame}

    \begin{frame}
        \frametitle{blub}
    %{\setlength{\tabcolsep}{1mm} 
    %\begin{tabu} to \linewidth{@{}XXXXX@{}}
        \begin{tabularx}{\linewidth}{@{\extracolsep{\fill}}cccccc}
    \footnotesize Input&
    \footnotesize Random Forest&
    \footnotesize SVM on SP&
    \footnotesize CRF&
    \footnotesize Ground Truth\\

    \includegraphics[width=.17\textwidth]{images/00845_image.png}&%
    \includegraphics[width=.17\linewidth]{images/00845_pixel.png}&%
    \includegraphics[width=.17\linewidth]{images/00845_svm.png}&%
    \includegraphics[width=.17\linewidth]{images/00845_ssvm.png}&%
    \includegraphics[width=.17\linewidth]{images/00845_gt.png}\\

    \includegraphics[width=.17\textwidth]{images/00781_image.png}&%
    \includegraphics[width=.17\linewidth]{images/00781_pixel.png}&%
    \includegraphics[width=.17\linewidth]{images/00781_svm.png}&%
    \includegraphics[width=.17\linewidth]{images/00781_ssvm.png}&%
    \includegraphics[width=.17\linewidth]{images/00781_gt.png}\\

    \includegraphics[width=.17\textwidth]{images/01331_image.png}&%
    \includegraphics[width=.17\linewidth]{images/01331_pixel.png}&%
    \includegraphics[width=.17\linewidth]{images/01331_svm.png}&%
    \includegraphics[width=.17\linewidth]{images/01331_ssvm.png}&%
    \includegraphics[width=.17\linewidth]{images/01331_gt.png}\\

    %\includegraphics[width=.15\textwidth]{images/00118_image.png}&%
    %\includegraphics[width=.15\linewidth]{images/00118_pixel.png}&%
    %\includegraphics[width=.15\linewidth]{images/00118_svm.png}&%
    %\includegraphics[width=.15\linewidth]{images/00118_ssvm.png}&%
    %\includegraphics[width=.15\linewidth]{images/00118_gt.png}

    \end{tabularx}
    %\multicolumn{5}{@{}c@{}}{Input Image}\vspace{3mm}\\

    %\multicolumn{5}{@{}c@{}}{Random Forest Prediction}\vspace{3mm}\\

    %\multicolumn{5}{@{}c@{}}{Superpixel Voting}\vspace{3mm}\\

    %\multicolumn{5}{@{}c@{}}{Support Vector Machine on Superpixels}\vspace{3mm}\\

    %\multicolumn{5}{@{}c@{}}{Conditional Random Field on Superpixels}\vspace{3mm}\\

    %\multicolumn{5}{@{}c@{}}{Ground Truth}\vspace{3mm}\\
    %\multicolumn{5}{@{}c@{}}{%
        %\includegraphics[]{legend.pdf}
    %}
    %\end{tabu}
%}
        % qualitative results
    \end{frame}


    \begin{frame}
        % Learned potentials
    \end{frame}


    \section*{Structured Machine Learning in Python}
    \begin{frame}
        \frametitle{PyStruct - simple structured prediction}
        \begin{columns}
            \column{.6\textwidth}
                \begin{itemize}
                    \item Implements common solvers: perceptron, cutting plane
                        ssvm, latent ssvm, subgradient, block-coordinate
                        Frank-Wolfe.
                    \item<2-> Implements powerful models for many use cases.
                    \item<3-> Interface to popular inference libraries: LibDAI, OpenGM, QPBO, AD3.
                    \item<4-> Lots of documentation and examples.
                \end{itemize}
            \column{.4\textwidth}
            \begin{center}
                \includegraphics[width=.6\linewidth]{images/logo_pystruct}
            \end{center}
        \end{columns}
    \end{frame}

    \begin{frame}
        \begin{center}
            \includegraphics[width=\linewidth]{images/example_gallery}\\
            http://pystruct.github.io
        \end{center}
    \end{frame}

    \begin{frame}
        \frametitle{Take Home}
        \begin{itemize}
            \item High level semantic unaries.
            \item Learn everything.
            \item Exact learning of CRF possible.
        \end{itemize}
    \end{frame}
    
    \section*{Thank you.}
    \begin{frame}
        \begin{center}
            Thank you for your attention.

            Questions?
        \end{center}
    \end{frame}
    \begin{frame}<beamer:0>
        \bibliographystyle{plainnat}
        \bibliography{paper}
    \end{frame}

	\end{document}
